\begingroup
\scriptsize
\setlength{\tabcolsep}{1.5pt}
\begin{longtable}{@{}>{\RaggedRight\arraybackslash}p{0.065\textwidth}>{\RaggedRight\arraybackslash}p{0.095\textwidth}>{\RaggedRight\arraybackslash}p{0.125\textwidth}>{\RaggedRight\arraybackslash}p{0.145\textwidth}>{\RaggedRight\arraybackslash}p{0.12\textwidth}>{\RaggedRight\arraybackslash}p{0.12\textwidth}>{\RaggedRight\arraybackslash}p{0.13\textwidth}>{\RaggedRight\arraybackslash}p{0.13\textwidth}@{}}
\caption{Progressive evidence chain from information characterization to evidence-calibrated decision realization.}
\label{tab:progressive-evidence}\\
\toprule
Part & Stage & Scientific question & Key evidence & Effect & 95\% CI / domains & Narrative conclusion & Control boundary \\
\midrule
\endfirsthead
\multicolumn{8}{@{}l}{\tablename~\thetable{} (continued)} \\
\toprule
Part & Stage & Scientific question & Key evidence & Effect & 95\% CI / domains & Narrative conclusion & Control boundary \\
\midrule
\endhead
\midrule
\multicolumn{8}{r@{}}{Continued on next page} \\
\endfoot
\bottomrule
\endlastfoot
Part I & I1 Spatial enrichment & Does spatial internal morphology enrich residual-capacity infor\allowbreak{}mation? & Matched scalar and surface references versus selected B-family spatial field & MAE 0.18920 to 0.12849; reduction 0.06072 (32.1\%) & familywise 95\% CI [0.00664, 0.15364]; favorable 5/\allowbreak{}6 domains & Spatial morphology adds CAI-relevant infor\allowbreak{}mation beyond lower-dimensional inputs. & The 0.08964 independent field estimate is sensitivity-only, not the matched confirmatory path. \\
 & I2 Sparse recoverability & How much registered full-field infor\allowbreak{}mation survives sparse observation? & Selected 25\% bilinear sparse field versus surface and full-field references & 89.9\% gain retained; sparse-to-full MAE gap 0.00602 & gap 95\% CI [0.00173, 0.01083]; favorable 5/\allowbreak{}6 versus surface and 0/\allowbreak{}6 versus full & Sparse observations preserve most of the registered spatial gain. & Native-raster fraction is not physical scanner time and sparse is not equivalent to full field. \\
 & I3 Spatial heterogeneity & Is task-relevant measurement opportunity spatially heterogeneous? & One-shot mechanical oracle versus uniform and recon\allowbreak{}struction oracles & CAI-AUEBC improvements 0.00391 and 0.00373; 56.8\% sequential headroom retained & 95\% CIs [0.00280, 0.00502] and [0.00284, 0.00469]; favorable 6/\allowbreak{}6 & Specimen-specific spatial acquisition headroom exists retrospectively. & All oracle rows use unavailable counterfactual outcomes and are non-deployable. \\
 & I4 Objective conditioning & Do downstream objectives induce distinct measurement priorities? & Mechanical and registered normalized-RGB-MSE recon\allowbreak{}struction oracles plus learned global masks & CAI contrast 0.04862; image-MSE contrast 5.503e-4; learned indicator 0 & 95\% CIs [0.04527, 0.05205] and [5.006e-4, 6.063e-4] & Retrospective priorities are objective-conditioned. & Learned global masks do not reproduce the oracle separation or establish deployment value. \\
 & I5 State conditioning & Does measurement value evolve with inspection state? & Strict-OOF retrospective teacher from initial state to 18.75\% checkpoint & 70.4\% best-action turnover; rank 0.405; top-5 overlap 0.307; opportunity 0.00531 & descriptive across 276 specimens and 6 held-out domains & A fixed initial ranking omits material state-conditioned value evolution. & Teacher evolution does not show that the deployable scorer observes the change. \\
 & I6 Predictor conditioning & How does downstream predictor choice bound the task-value definition? & Ridge-Huber and Ridge-shallow-MLP strict-OOF value agreement & Spearman 0.762 versus 0.116; full-state MAE 0.08964, 0.08618, 0.15067 & rank CIs [0.699, 0.821] and [0.069, 0.164]; 6 held-out domains & Value is stable between comparably performing low-complexity predictors. & The shallow MLP is substantially less accurate; equally accurate structurally distinct predictors remain unresolved. \\
\addlinespace[2pt]
Part II & II1 Static reference & What reference motivates state-conditioned valuation? & Static legal-state scorer with global and random exact-budget controls & static Spearman -0.0196; set regret 0.08171 versus 0.07993 and 0.07978 & rank 95\% CI [-0.0591, 0.0195]; favorable 3/\allowbreak{}6 domains & The static reference motivates conditional valuation. & This tested repre\allowbreak{}sentation does not imply infor\allowbreak{}mation-theoretic impossibility. \\
 & II2 Dynamic valuation & Does state-conditioned scoring capture incremental dependence? & Dynamic real-state scorer versus static scorer at 18.75\% & dynamic minus static regret -0.001260 & 95\% CI [-0.002123, -0.000444]; favorable 5/\allowbreak{}6 domains & Dynamic valuation improves over the static reference. & The comparison alone does not identify which state-infor\allowbreak{}mation source drives the gain. \\
 & II3 Infor\allowbreak{}mation-source attribution & Which legal-state infor\allowbreak{}mation licenses content attribution? & Real content versus initial, acquired-position/\allowbreak{}history, recon\allowbreak{}struction, and shuffled controls & real change -0.000731; real-minus-history 0.01740; real-minus-recon\allowbreak{}struction 0.03419; real-minus-shuffled 2.328e-4 & all matched-control CIs exclude 0; real favorable 1/\allowbreak{}6 for each source control & Real measurements change prediction, while matched controls localize the decision signal. & Acquired-position/\allowbreak{}history is not geometry-only; recon\allowbreak{}struction is registered normalized-RGB-MSE; shuffled remains the dynamic boundary. \\
 & II4 Valuation/\allowbreak{}planning decomposition & Where does controlled component headroom reside? & Privileged valuation and bounded planning substitutions & valuation 4.979e-5; learned planning 3.164e-6; true-value planning 1.117e-4 & 95\% CIs [1.845e-5, 8.105e-5], [1.411e-7, 6.250e-6], [9.898e-5, 1.248e-4] & Valuation and planning contribute distinguishable retrospective headroom. & Substitutions are diagnostic and non-deployable; no repre\allowbreak{}sentation bottleneck is claimed. \\
 & II5 Bounded set realization & How much decision-level headroom remains in the reachable pool? & Greedy and beam-width-four selection versus retrospective joint near-oracle & greedy regret 1.207e-4; beam-4 point regret 1.127e-4 & 95\% CIs [1.033e-4, 1.386e-4] and [9.485e-5, 1.308e-4] & A positive two-action set-realization gap remains. & The near-oracle is non-deployable and does not establish arbitrary-horizon regret. \\
 & II6 Deployment calibration & How does frozen end-to-end stress testing calibrate readiness? & Feedback/\allowbreak{}no-feedback stress test and strongest deployable outer reference & no-feedback advantage 1.496e-5; residual deployable gap 6.114e-5 & feedback-benefit CI [-1.944e-5, -1.064e-5]; baseline-minus-learned CI [-8.461e-5, -3.777e-5] & The frozen endpoint quantifies the remaining repre\allowbreak{}sentation-to-decision gap. & The no-feedback reference is retained and the learned endpoint is not performance-superior. \\
\end{longtable}
\endgroup
